\section*{Player Agency}\stepcounter{section}\phantomsection\addcontentsline{toc}{section}{Player Agency}
{\entryfont Throughout the debate, the player characters should be free to contribute whenever they see an opportunity. They may challenge a speaker, support one of the positions already raised, combine elements from several viewpoints, or introduce an entirely different perspective. Encourage the discussion to develop naturally rather than treating the characters' involvement as a separate stage that begins only after the council has finished speaking.

The council should respond to the characters as part of the ongoing debate. Speakers may question their suggestions, ask them to elaborate, or refer back to arguments raised earlier in the discussion. This allows the characters' views to become woven into the council's deliberations rather than presented as isolated statements after the fact.

If the players remain hesitant or none of the characters has entered the discussion, a council member may eventually turn toward them and ask:}
\begin{DndReadAloud}
	"If the portals never return, what should become of the Caelthir?"
\end{DndReadAloud}
{\noindent\entryfont This question may also be introduced organically during the preceding debate whenever it provides a natural opening for the characters to participate.

Persuasion, Insight, History, or similar checks may be appropriate when a character attempts to influence a particularly hesitant speaker, understand the reasoning behind an argument, or support their position with relevant knowledge. Such checks should complement the roleplaying rather than determine which viewpoint is considered correct.

The discussion does not need to reach a unanimous conclusion. Its purpose is to establish how the characters themselves view the changing circumstances of Ty Coeden and to give the players a personal stake in the decisions that follow.}